\cotefile{60}
\scanpage{14}\mspage{- 9 -}\hfill n° 57

\uline{Espace des sous-groupes fermés}, (n° 384)

C'est adopté, et ça ira en \S{} 5 du chap. VIII d'Intégration. On vide l'espace uniforme des parties fermées et la prop. 3 (passent en exerc., avec appel d'exerc.).

Godement a fait une nouvelle rédaction en Congrès, qui a été relue, et améliorée quant aux virgules ; le rédacteur définitif a noté celles-ci.

\uline{Ce qui devrait être le chapitre des groupes de Lie} (n° 385 et 386).

Le n° 385 (``Distributions ponctuelles sur une variété''), sera incorporé dans le Fasc. Rés. des Variétés, après avoir été considérablement allégé, On n'y parlera pas du cas algébrique (ce qui simplifie en bas de la p. 3, le corps étant infini). Naturellement, ça doit être refondu avec les jets. On dira ``distributions à support fini''. Alléger les sorites (p. ex. celui des produits, n° 8 ; pourquoi pas les limites inductives et projectives ??).

L'énorme n° 386 (``Groupes formels et hyperalgèbres'') n'est pas du tout ce que Bourbaki voulait. D'une part le lien entre groupes analytiques et algèbres de Lie n'est fait que dans un mythique chapitre ultérieur. D'autre part on n'y démontre pas de résultats substantiels sur les groupes formels en caract. p. Ainsi Bourbaki vide toute la rédaction, et, comme Edith Piaf, ``on repart à zéro''. Cependant, le \S{} 4 prévu par Cartier (correspondance entre groupes formels et algèbres de Lie en caract. 0, formule de Hausdorff), nettement plus substantiel que le reste, devra être rédigé pour usage ulterieur.

Là dessus Bourbaki s'aperçoit qu'il dispose d'une fort bonne rédaction de Koszul (``Variétés de groupes'', n° 340), sur laquelle on a intelligemment travaillé (voir Tribu n° 51). Ça doit facilement s'articuler avec le Fasc. Rés. Variétés. Bruhat devra rédiger sur cette base, sans chercher à faire d'originalité. Sinon le Livre de Lie est fichu !

\begin{center}
-:-:-:-
\end{center}

\begin{center}
\uline{Séminaire BOURBAKI}
\end{center}

Décembre 62 :

\begin{tabular}{@{}l@{\quad}p{.75\linewidth}@{}}
CARTIER & - Fluctuations dans les suites de variables aléatoires indépendantes\\
DEJEAN & - Papier G\handwritten{\r{a}}rding à la S. M. F.\\
DIEUDONNÉ & - Travail de Kostant (sur Amitsur - ... etc)
\end{tabular}
% archive sheet number: 13
\scanpage{15}
\[
L_{DE}(t)=L_E\bigl(p^{-(n\unsure{-|+}\rho)}t\bigr)
\]
d'où, portant \unsure{dans (2)|du cas} et remplaçant $t$ par $t^{-1}$ :
\[
\boxed{L_E\left(\frac{1}{p^{n+\rho}t}\right)=\varepsilon(E)\left(-p^{\frac{n+\rho}{2}}t\right)^{\chi(E)}L_E(t)}\tag*{(3)}
\]

Posant
\[
\boxed{\xi_E(t)=t^{\chi(E)/2}L_E(t)}
\]
il vient l'équation fonctionnelle
\[
\boxed{\xi\left(\frac{1}{p^{n+\rho}t}\right)=\varepsilon'(E)\,\xi(t)}
\]
\[
\begin{aligned}
\varepsilon'(E)=\pm1,\qquad \varepsilon'(E)&=(-1)^{\chi(E)}\prod_i\operatorname{signe}\operatorname{dét} f_{H^i(\bar X,\bar F)}\\
&=(-1)^{\chi(E)}\operatorname{signe}\operatorname{det} f_{H^*(\bar X,\bar F)}
\end{aligned}
\]

\deleted{Si on \unsure{veut} l'équation fonctionnelle avec $\varepsilon(E)$ au lieu de} Notons d'ailleurs que $H^i(\bar X,\bar E)$ et $H^{2n-i}(\bar X,\check{\bar E})$ \inserted{$\simeq H^{2n-i}(\bar X,\bar E(\rho))=H^{2n-i}(\bar X,\bar E)(\rho)$} étant ``\unsure{duals|duaux|dual} \unsure{à|:} valeurs \unsure{dans|au} $\mathbb{Q}_{\ell}(n)$'', et l'opération de $f$ sur \inserted{$\mathbb{Q}_{\ell}(\unsure{1|n})$} \unsure{--|en|étant} de signe positif, on trouve que $\varepsilon_i(E)=\varepsilon_{2n-i}(E)$, et par suite
\[
\boxed{\varepsilon(E)=\operatorname{Signe}\operatorname{det} f_{H^n(\bar X,\bar E)}}.
\]
Pour la raison analogue, $(-1)^{\chi(E)}=(-1)^{b_n}$. Donc
\[
\varepsilon'(E)=(-1)^{b_n(E)}\operatorname{signe}\operatorname{dét}f_{H^n(X,E)}
\]
% archive sheet number: 14
\scanpage{16}\hfill Rédaction n° 398

\begin{center}
COMPLÉMENTS AU CHAPITRE VII D'ALGÈBRE

DUALITÉ DES MODULES DE LONGUEUR FINIE SUR UN ANNEAU PRINCIPAL

-:-:-:-

CONTRE-RÉDACTION
\end{center}

[On se refuse à parler de modules injectifs, et d'autre part on se borne aux modules de longueur finie. Cela donne une rédaction minimale, (dans l'ensemble des rédactions $\ne\varnothing$)].

\uline{Notations}.

$A$ désigne un anneau principal, $K$ son corps des fractions. On suppose $K\ne A$ : i.e., $A$ possède au moins un élément extrémal.

Si $E$ est un $A$-module, on pose $E^*=\operatorname{Hom}(E,K/A)$ ; c'est un $A$-module. A tout homomorphisme $E\longrightarrow F$ correspond un homomorphisme $F^*\longrightarrow E^*$. Si $x\in E$, $x^*\in E^*$, on pose $x^*(x)=(x,x^*)$. On définit ainsi une application $A$-bilinéaire de $E\times E^*$ dans $K/A$, d'où un homomorphisme :
\[
i_E:E\longrightarrow E^{**}.
\]

PROPOSITION 1.- \uline{Si} $E$ \uline{est un} $A$-\uline{module de longueur finie}, $E^*$ \uline{est isomorphe}, (non canoniquement en général), \uline{à} $E$, \uline{et l'application} $i_E:E\longrightarrow E^{**}$ \uline{est un isomorphisme}.

Vu les théorèmes de structure, on peut se borner au cas où $E=A/tA$, avec $t\ne0$. Dans ce cas $E^*$ s'identifie au sous-module de $K/A$ formé des $x$ tels que $tx=0$, i.e. au sous-module $t^{-1}A/A$, isomorphe à $A/tA$. D'où l'isomorphie de $E$ et de $E^*$. Enfin, lorsque $E=A/tA$, la détermination précédente montre que $E^{**}$ s'identifie à $A/tA$ et l'on vérifie tout de suite que, dans ce cas, $i_E$ est l'application identique. D'où la deuxième partie de la proposition.
% archive sheet number: 15
\scanpage{17}
\uline{Si} $n+\rho$ \uline{est impair}, alors $b_n=\dim\underbrace{H^{\unsure{n+\rho|n}}(\unsure{\overline{\bar X,\bar E}|\bar X,\bar E})}_{\text{poids }n+\rho}$ est pair, donc $\deleted{\unsure{b_n|b}}(-1)^{b_n}=1$. De plus, un raisonnement du cours de Serre utilisant une forme alternée non dégénérée invariante montre que dans ce cas $\operatorname{signe}\operatorname{dét}f_{H^n(\unsure{X|\bar X},\bar E)}$ est \unsure{également|égal à} $+1$. Donc on trouve alors $\varepsilon'(E)=1$.

Si $n+\rho$ est pair, on ne peut déterminer le signe de façon absolue. On peut \unsure{seult|seul\textsuperscript{t}} \unsure{obs.|obs} \inserted{(indép. de la parité)} \deleted{que} si $\varepsilon_1$ $(\varepsilon_{-1})$ est la \deleted{\unsure{par dimension|dimension|par dimension de}} \inserted{multiplicité} de \deleted{\unsure{l'esp.|l'espace}} la valeur propre $1$ $(-1)$ de $f$ sur $\unsure{H^n|H}(\bar X,\bar F)$, on a $(-1)^{b_n}=(-1)^{\varepsilon_1+\varepsilon_{-1}},\quad\varepsilon(E)=(-1)^{\varepsilon_{-1}}$, donc
\[
\boxed{\varepsilon'(E)=(-1)^{\varepsilon_1}}.
\]

Lorsque l'on n'est pas dans le cas favorable $X$ lisse, $E$ \unsure{ct tordu|ct tordus|de torsion}, $\deleted{\unsure{E|F}}\check E$ \uline{globalement} isomorphe à $E(\rho)$, alors $D(E)$ ne s'exprime plus de façon aussi simple en fonction de $E$, et l'équation fonctionnelle (1) ne s'exprime plus \unsure{néc.|nic.} en termes de la seule fonction $L_E$ (mais elle relie \uline{deux} fonctions $L_{DE},L_E$]
% archive sheet number: 16
\scanpage{18}\mspage{- 8 -}
\hfill n° 57

\diagram{\begin{tabular}{|l@{\quad}c@{\quad}p{.68\linewidth}|}
\hline
DIEUDONNÉ & : & Rédactions définitives\newline Systèmes de Racines (Juin 63)\\ \hline
DIXMIER & : & Théorie fine des algèbres normées (mars 63)\\ \hline
GODEMENT & : & Rapport d'Analyse dirigée vers les groupes semi-simples (Juin 63)\\ \hline
KOSZUL & : & Géométrie Riemannienne locale et globale (Juin 63)\\ \hline
LANG & : & Rapport sur Mlle Bastiani \handwritten{?}\\ \hline
SAMUEL & : & La Tribu\\ \hline
SCHWARTZ & : & Rapport sur les Probas (?)\\ \hline
SERRE & : & Semi-simples complexes (Juin 63)\newline Papier sur l'approximation forte dans SL\\ \hline
TATE & : & Dimension (mars 63),\\ \hline
WEIL & : & Rien\\ \hline
\end{tabular}}

\begin{center}-:-:-\end{center}

\begin{center}\uline{DÉCISIONS LOCALES DIVERSES}\end{center}

\uline{Errata}.

Pour Ens. III, p. 26 : remplacer `` flirtant à droite et à gauche '' par `` flirtant à gauche et à droite ''.

Pour Alg. Commut. I, p. 55, ne pas tromper le lecteur quant aux intentions de Bourbaki, et rédiger ainsi la footnote : `` En attendant l'éventuelle publication d'une partie de ce Traité consacrée à l'Algèbre Homologique, le lecteur pourra consulter ... ''.
% archive sheet number: 17
\scanpage{19}
$X$ lisse \unsure{purement|purem\textsuperscript{t}} de dim $n$

$E$ \unsure{$c^{t}$|ct|et} tordu purement de poids $\rho$

$X\xrightarrow{f}\operatorname{Spec}\mathbb F_p$ \unsure{comp. actifiable|compactifiable}

\[
\text{\uline{Dualité}}\qquad \unsure{\mathbb D|D}\,\mathbb R f_!(E)=\mathbb R f_*(DE)
\]
Or par hyp. sur $E$
\[
DE=\check E(n)[2n]\simeq E(\rho+\deleted{\unsure{2|2n}}n)[2n]
\]
D'où
\[
D\,\mathbb R f_!(E)=\bigl(\mathbb R f_*(E)\bigr)(\rho+n)[2n]
\]
d'où en prenant les classes, et écrivant
\[
\begin{aligned}
\operatorname{cl}D(\mathbb R f_!(E))&=\bigl(\operatorname{cl}\mathbb R \unsure{f_*|f}(E)\bigr)\cdot\deleted{\unsure{\xi^{\rho+n}|p^{\rho+n}}}(\rho+n)\\
&=\operatorname{cl}\bigl(\mathbb R f_!(E)\bigr)(\rho+n)+\operatorname{cl}\mathbb R_\infty f(E)(\unsure{\rho\,n|\rho+n})
\end{aligned}
\]
d'où, en prenant les fonctions $L$
\[
\deleted{L_{DE}}(-t)^{-\chi(E)}\deleted{\unsure{\delta|d}}\delta(E)L_E(\tfrac{1}{t})=L_E\bigl(p^{-(\rho+n)}t\bigr)L_E^{\unsure{(\infty|\infty}}\bigl(p^{-(\rho+n)}\unsure{t|1}
\]
ou encore
% archive sheet number: 18
\scanpage{20}\mspage{6 - 44 -}\hfill n° 392

est principal. Inversement, pour un nombre premier $p$, considérer le sous-anneau $C=Z+Bp$ de $B$ ; on a $\underline f=\underline f_{C/B}=Bp$, et $C/\underline f$ est isomorphe à $\underline Z/p\underline Z=\underline F_p$ ; observer enfin que $C/\underline f$ est un $C$-module isomorphe à $C^*/B^*$, donc à $B/C$, et en conclure que $B$ est nécessairement un $\underline Z$-module de dimension 2).

6) Les hypothèses sur $A$, $B$, $C$ étant celles de l'exerc. 2, soient $\underline p$ un idéal premier de $A$ de hauteur 1, $\underline P$ un idéal premier de $C$ au-dessus de $\underline p$. Montrer que les conditions suivantes sont équivalentes :

a) $\underline P_{\underline p}$ ne contient pas $(\underline D_{C/A})_{\underline p}$.

b) $\underline P_{\underline p}$ est une des composantes primaires de $\underline p C_{\underline p}$ dans $C_{\underline p}$, et le corps des fractions de $C/\underline P$ est une extension séparable de celui de $A/\underline p$.

c) Il existe un seul idéal premier $\underline P'$ de $B$ au-dessus de $\underline P$, $\underline P'$ est non ramifié sur $\underline p$, et les corps des fractions de $C/\underline P$ et de $B/\underline P'$ sont identiques.

(Se ramener au cas où $A$ est un anneau de valuation discrète. Pour voir que a) entraîne b), utiliser l'exerc. 3 b), l'exerc. 1 b) pour évaluer la longueur de $C/\underline Q$, où $\underline Q$ est la composante $\underline P$-primaire de $C\underline p$, et enfin le chap. V, \S{} 1, exerc. 16. Utiliser encore l'exerc. 1 b) pour voir que b) entraîne c). Enfin, pour montrer que c) implique a), considérer les complétés des anneaux semi-locaux $B$ et $C$).

7)\quad a) Sous les hypothèses de la prop. \quad, montrer que si les $x^i$, $(0\leq i\leq n-1)$, forment une base du $A$-module $B$, il en est de même des $y^i$, $(0\leq i\leq n-1)$, pour tout $y\in B$ suffisamment voisin de $x$.

b) Soient $A$ un anneau de valuation discrète, $K$ son corps des fractions, $B$ la fermeture algébrique de $A$ dans une extension finie séparable $L$ de $K$. Soit $\underline P$ un idéal premier de $B$ tel que le corps des fractions de $B/\underline P$ soit une extension séparable du corps résiduel $\bar K$ de $A$. Montrer qu'il existe $x\in B$ tel que $L=K(x)$ et tel que le conducteur $f_{B/A[x]}$ de $B$ dans $A[x]$ soit étranger à $\underline P$.

c) Les hypothèses sur $A$ et $B$ étant celles de b), supposons qu'il existe $n=[L:K]$ idéaux premiers distincts de $B$ au-dessus de l'idéal maximal $\underline p$ de $A$, (on dit alors que $B$ est ``complètement décomposé''). Montrer que, pour qu'il existe $x\in B$ tel que $B=A[x]$, il faut et il suffit que l'on ait $n\leq\operatorname{Card}(\bar K)$, (utiliser le lemme de Hensel et le th. d'approximation du chap. VI, \S{} 7, n° \quad).
% archive sheet number: 19
\scanpage{21}
\margin{en admettant la résolution des singularités ... et les conjectures de Weil\par\rule{.32em}{.3pt}\ \rule{.65em}{.3pt}\ \rule{.43em}{.3pt}}
\[
\boxed{L_E\left(\frac{1}{p^{n+\rho}t}\right)=A(t)\,\delta(E)\,(-t)^{\chi(E)}\,L_E(t)}
\]
\[
\left\{\begin{aligned}
A(t)&=L_E^\infty\left(\deleted{p}\frac{1}{p^{n+\rho}t}\right)^{-1}\\
&\deleted{\unsure{=|\simeq}}\\
\chi(E)&=\sum(-1)^i\operatorname{rang}H^i_!(\bar X,\bar F)\\
\delta(E)&=\prod_i\deleted{\unsure{(-1)^i|\varepsilon_i}}\bigl(\operatorname{dét}\operatorname{frob}_{H^i_!(\bar X,\bar F)}\bigr)^{(-1)^i}
\end{aligned}\right.
\]

\uline{NB}\quad $L_E^\infty(t)=\inserted{=L^\infty_{\mathbb R_\infty f_*(E)}(t)=}L_{\mathbb R i_*(E)|Y}(t)$ est de \deleted{\unsure{niveau|niveau $\rho$}} \inserted{type} $\leq\tau+(n-1)$ ($\tau$ étant le type de $E$) et de poids $\leq\rho+2n$

\uline{Si $X$ complet}
\[
A(t)=1
\]

On peut aussi écrire
\[
\chi(E)=\varepsilon(\mathbb R f_!(E)),\quad \deleted{\unsure{\varepsilon|=}}\delta(E)=\deleted{\unsure{F|f}\,\operatorname{dét}\unsure{\operatorname{frob}|f}}\delta(\mathbb R f_!(E))
\]
\[
\varepsilon(\mathbb R f_!(E))=\varepsilon(D\mathbb R f_!(E))\deleted{\unsure{=|\varepsilon}},\quad D\mathbb R f_!(E)=\mathbb R\unsure{f|f_*}(DE)=(\mathbb R f_*(E))(\rho+\deleted{\unsure{2|2n}}n)\inserted{[2n]}
\]
\[
\deleted{\operatorname{det}}\delta(\mathbb R f_!(E))=\delta(D(\mathbb R f_!(E))^{-1}
\]
d'où \deleted{d'où}
\[
\deleted{\unsure{\lbrace|f}}\left\{\begin{aligned}
\chi(E)&=\sum(-1)^i\operatorname{rang}\mathbb R^i f_*(E)\\
\delta(E)&=p^{+(\rho+n)\chi(E)}\delta(\mathbb R f_*(E))^{-1}=p^{(n+\rho)\chi(E)}\prod_i\deleted{\unsure{(-1)^i|\delta_i}}\bigl(\operatorname{det} f_{H^i(\bar X,\bar F)}\bigr)^{(-1)^{i+1}}
\end{aligned}\right.
\]

[Comparant, lorsque $X$ est complet, les deux valeurs obtenues de $\delta(E)$, on trouve dans ce cas
\[
\delta(E)^2=p^{(n+\rho)\chi(E)}\qquad (X\text{ complet})
\]
ce qui est bien compatible avec les conjectures de Weil sur les val. absolues des valeurs propres.... ]
% archive sheet number: 20
\scanpage{22}\mspage{6 - 43 -}\hfill n° 392

e) Pour que $C^*$ soit pseudo-inversible, il faut et il suffit que pour tout idéal premier $\underline p$ de hauteur 1 de $A$, on ait $(\underline D_{C/A})_{\underline p}:(\underline D_{C/A})_{\underline p}=C_{\underline p}$.

4)\quad a) Soient $C$ un anneau noethérien intègre, $\unsure{\underline q|\deleted{\underline q}}$ un idéal primaire tel que $\underline q:\underline q=C$, $\underline p$ l'idéal premier correspondant ; on a alors $\underline q:\underline p\subset C$. Montrer que pour que l'idéal primaire $\underline q$ soit irréductible, il faut et il suffit que $(\underline q:\underline p)/\underline q$ soit un $C$-module simple. (Se ramener au cas où $C$ est un anneau artinien primaire avec $\underline q=0$, $\underline p$ étant donc l'unique idéal premier de $C$ ; montrer que $0:\underline p$ est le socle de $C$, en notant que c'est un $(A/\underline p)$-module, et pour que $0$ soit irréductible, il faut et il suffit que le socle de $C$ soit simple).

b) On suppose que les anneaux $A$, $B$, $C$ vérifient les hypothèses de l'exerc. 2. Montrer que, pour que l'idéal fractionnaire $C^*$ de $C$ soit pseudo-inversible, il faut et il suffit que pour tout idéal premier $\underline p$ de hauteur 1 dans $A$, les idéaux primaires de la décomposition primaire réduite, (unique), de $\underline p C_{\underline p}$ dans $C_{\underline p}$, soient irréductibles. (Se ramener au cas où $A$ est un anneau de valuation discrète d'uniformisante $\pi$ ; si $\underline m_i$, $(1\leq i\leq r)$, sont les idéaux maximaux de $C$, pour que $C^*$ soit inversible, il faut et il suffit que les $C$-modules $C^*/\underline m_i C^*$ soient simples, (cf. exerc. 1 b) ou, ce qui revient au même, que les $C$-modules $(\pi C)^*/\underline m_i(\pi C)^*$ soient simples ; utilisant la relation $(\underline a\underline b)^*=\underline a^*:\underline b$, en déduire que cette condition équivaut à dire que $(\pi C:\underline m_i)/\pi C$ est simple pour tout $i$ ; enfin, si $\underline q_i$ est la composante $\underline m_i$-primaire de $C$, observer que $\pi C:\underline m_i=\underline q_i:\underline m_i$ et utiliser a)).

c) Afin que, pour tout idéal premier $\underline p$ de hauteur 1 de $A$, les $C_{\underline p}$-modules $C^*_{\underline p}/C_{\underline p}$ et $C_{\underline p}/(\underline D_{C/A})_{\underline p}$ soient isomorphes, il faut et il suffit que l'idéal fractionnaire $C^*$ soit pseudo-inversible. (Se ramener au cas où $A$ est un anneau de valuation discrète ; observer que si $\underline m_i$, $(1\leq i\leq r)$, sont les idéaux maximaux de $C$, l'hypothèse que $C^*/C$ est isomorphe à $C/\underline D_{C/A}$ entraîne que $C^*/\underline m_i C^*$ est un $C$-module simple pour tout $i$).

d) Montrer que pour que l'on ait $C_{\underline p}=B_{\underline p}$ pour tout idéal premier $\underline p$ de hauteur 1 de $A$, il faut et il suffit que tout idéal primaire $\ne 0$ de $C$ soit pseudo-inversible. (Pour tout idéal premier $\underline p\ne0$ de $C$, considérer l'idéal $\underline p^2$).

5) Soit $K$ une extension algébrique finie de $Q$, $B$ la fermeture intégrale de $Z$ dans $K$. Afin que, pour tout sous-anneau $C$ de $B$ ayant même corps des fractions que $B$, $C^*$ soit un idéal fractionnaire inversible de $C$, il faut et il suffit que $[K:Q]\leq2$. (Si $[K:Q]=2$, montrer, en utilisant la prop. \quad, que $\underline D_{C/Z}$
% archive sheet number: 21
\scanpage{23}
\fbox{Dans le cas d'un \uline{motif} $E$,}

On va exprimer $\delta(E)$ et $\chi(E)$ en termes des \deleted{\illegible{1}} composantes virtuelles des $\mathbb R f_!(E)$ (utilisant conj. de Weil), on trouve
\[
\left\{\begin{aligned}
\chi(E)&=\sum(-1)^i\beta_i(E)\\
\delta(E)&=\prod_i\delta\bigl(h^i_!(E)\bigr)^{(-1)^i}
\end{aligned}\right.
\]
\margin{$\beta_i(E)$ les \unsure{nbs|nb\textsuperscript{s}} de Betti virtuels à coeff. dans $E$}

Or $h^i_!(E)$ donne $\beta_i$ valeurs propres de val. absolue $p^{i/2}$, donc
\[
\deleted{\unsure{\delta(\beta_i)=|\delta(\beta)=}}\delta\bigl(h^i_!(E)\bigr)=\varepsilon_i\,p^{\frac{i\beta_i}{2}}
\]
\margin{(NB. $i\beta_i/2$ tjrs entier)}

où $\deleted{\unsure{\delta|\varepsilon}}\varepsilon_i$ est défini comme un signe de déterminant virtuel, et $\varepsilon_i$ est $\deleted{\unsure{=|\varepsilon}}+1$ si $i$ est impair. On trouve alors
\[
\left\{\begin{aligned}
\delta(E)&=\deleted{\unsure{\varepsilon(E)p^{\frac12\sum i\beta_i}|\varepsilon(E)p^{\sum i\beta_i}}}\varepsilon(E)\,p^{\frac12\sum_i(-1)^i\,i\,\beta_i(E)}\\
\varepsilon(E)&=\prod_i\varepsilon_i(E)
\end{aligned}\right.
\]
\margin{$\varepsilon(E)$ et les $\varepsilon_i(E)$ sont $\pm1$, $\varepsilon_i(E)=+1$ si $i$ impair.}
% archive sheet number: 22
\scanpage{24}\mspage{6 - 42 -}\hfill n° 392

2) Les hypothèses sur $A$ et $B$ étant les mêmes que dans l'exerc. 1, soit $C$ un anneau tel que $A\subset C\subset B$ et que le corps des fractions de $C$ soit égal à $L$ ; tout idéal $\ne0$ dans $C$ est alors un \uline{réseau} de $L$, (relativement à $A$).

a) Montrer que si $\underline a$, $\underline b$ sont deux idéaux de $C$ tels que $\underline a$ soit pseudo-inversible, on a $\chi(\underline a\underline b,C)=\chi(\underline a,C)+\chi(\underline b,C)$.

b) On dit que deux idéaux $\underline a$, $\underline b$ de $C$ sont \uline{pseudo-étrangers} si, pour tout idéal premier $\underline p$ de hauteur 1 dans $A$, $\underline a_{\underline p}$ et $\underline b_{\underline p}$ sont étrangers dans $C_{\underline p}$. Montrer que si $\underline a$ et $\underline b$ sont pseudo-étrangers et $\ne0$, on a $\chi(\underline a\underline b,C)=\chi(\underline a,C)+\chi(\underline b,C)$. (Se ramener au cas où $A$ est un anneau de valuation discrète, observer que l'on a $\unsure{\chi|\chi'}(\underline a,\underline a+\underline b)=\chi(\underline a\cap\underline b,\underline b)$ et utiliser la relation (\quad) du n° \quad du \S{} 4).

c) Déduire de b) que si $\underline a$ est un idéal $\ne0$ de $C$, et si $\underline q_i$, $(1\leq i\leq r)$, sont les idéaux primaires de la décomposition primaire réduite de $\underline a$, correspondant aux idéaux premiers non immergés, on a $\chi(\underline a,C)=\sum_i\chi(\underline q_i,C)$.

d) Soient $\underline q$ un idéal primaire de $C$, $\underline p$ l'idéal premier correspondant. Si $\underline p$ n'est pas de hauteur 1, on a $\chi(\underline q,C)=0$ ; dans le cas contraire, si $\underline p_{\unsure{o|0}}=\underline p\cap A$, on a $\chi(\underline q,C)=\operatorname{deg}_A(\underline q).\underline p_{\unsure{o|0}}$.

e) Soit $\underline a$ un idéal de $C$ ; pour que $\underline a$ soit pseudo-inversible, (exerc. 1 b)), il faut et il suffit que l'on ait $\chi(\underline a,C)=\chi(\underline a B,B)$. (Même méthode que dans l'exerc. 1 b)).

3) Les hypothèses sur $A$, $B$ et $C$ sont les mêmes que dans l'exerc. 2. On dit encore que le réseau $C^*$ est la \uline{codifférente} de $C$, et que $C:C^*$, qui est un idéal de $C$, est la \uline{différente} de $C$ sur $A$ ; on la note $\underline D_{C/A}$.

a) Montrer que le \uline{conducteur} $\underline f_{C/B}$ de $B$ dans $C$, (chap. V, \S{} 1, n° \quad), est égal à $B^*:C^*$, et que pour tout idéal oremier $\underline p$ de hauteur 1 dans $A$, on a $(\underline f_{C/B})_{\underline p}C^*_{\underline p}=B^*_{\underline p}$, (utiliser le fait que $B_{\underline p}$ est un anneau principal).

b) Montrer que l'on a $\underline D_{C/A}\subset\underline f_{C/B}\cap\underline D_{B/A}$, (utiliser a)).

c) Montrer que, pour que $C^*$, considéré comme idéal fractionnaire de $C$, soit pseudo-inversible, (exerc. 1 b)), il faut et il suffit que pour tout idéal premier $\underline p$ de hauteur 1 dans $A$, on ait $B_{\underline p}(\underline D_{C/A})_{\underline p}=(f_{C/B})_{\underline p}(\underline D_{B/A})_{\underline p}$, (utiliser a)).

d) Si $C^*$ est pseudo-inversible, tout idéal fractionnaire $\underline a\ne0$ de $C$ tel que $\underline a:\underline a=C$ est pseudo-inversible, (utiliser le fait que pour deux idéaux fractionnaires non nuls $\underline a$, $\underline b$ de $C$, on a $(\underline a\underline b)^*=\underline a^*:\underline b$).
% archive sheet number: 23
